\documentclass[sigconf]{cidr-2027}

\usepackage{microtype}
\usepackage{enumitem}
\usepackage{xspace}
\usepackage{booktabs}
\usepackage{tikz}
\usepackage{colortbl}

\setlist{nosep,leftmargin=1.25em}

\newtheorem{theorem}{Theorem}
\newtheorem{proposition}{Proposition}
\newtheorem{lemma}{Lemma}
\newtheorem{corollary}{Corollary}
\newtheorem{definition}{Definition}
\theoremstyle{remark}
\newtheorem{example}{Example}
\newtheorem{remark}{Remark}

\newcommand{\InputOrd}{\preceq_I}
\newcommand{\OutputOrd}{\preceq_O}
\newcommand{\InputGe}{\succeq_I}
\newcommand{\OutputGe}{\succeq_O}
% Output event emitting #2 after applying the \commitment sequence #1
% (\varepsilon = no \commitment).
\newcommand{\Out}[2]{\mathsf{output}_{#1}(#2)}
\newcommand{\commitment}{commitment\xspace}
\newcommand{\commitments}{commitments\xspace}
\newcommand{\Commitment}{Commitment\xspace}
\newcommand{\Commitments}{Commitments\xspace}
\newcommand{\LiftedOrd}{\trianglelefteq}
% Order vocabulary (role-based names; change the words here).
% Input order:  "j \inrel i" means i \InputOrd j (non-strict); inputs \ingrow;
%   a common \indesc of inputs \inrel each of them (articles: "an \indesc").
% Output order: "u \outrel o" means o \OutputOrd u; a common \outnoun of
%   outputs refines each of them.
\newcommand{\inrel}{expands\xspace}     % relation, 3rd person
\newcommand{\inrelbase}{expand\xspace}  % relation, base form
\newcommand{\inrelpart}{expanding\xspace}% relation, participle
\newcommand{\inrelpast}{expanded\xspace}% relation, past participle
\newcommand{\ingrows}{expands\xspace}        % growth, 3rd person
\newcommand{\ingrow}{expand\xspace}          % growth, base form
\newcommand{\inmore}{more expanded\xspace}    % informal comparative
\newcommand{\indesc}{expansion\xspace}      % noun
\newcommand{\indescs}{expansions\xspace}
\newcommand{\Indesc}{Expansion\xspace}
\newcommand{\outrel}{refines\xspace}         % relation, 3rd person
\newcommand{\outrelbase}{refine\xspace}      % relation, base form
\newcommand{\outrelpart}{refining\xspace}    % relation, participle
\newcommand{\outrelpast}{refined\xspace}     % relation, past participle
\newcommand{\outnoun}{refinement\xspace}     % noun
\newcommand{\outnouns}{refinements\xspace}
\setlength{\emergencystretch}{1.5em}
\newcommand{\Spec}{\mathcal{S}}
% The name of the specification model; optional suffix supports \spectype[s].
\newcommand{\spectype}[1][]{I/O specification#1\xspace}
\newcommand{\Spectype}[1][]{I/O specification#1\xspace}
\newcommand{\fiber}[1]{P(#1)}
\newcommand{\dog}{\texttt{dog}}
\newcommand{\catv}{\texttt{cat}}
\newcommand{\Subsumed}{\mathit{Subsumed}}
\newcommand{\Excluded}{\mathit{Excluded}}
\newcommand{\grayrow}[1]{\textcolor{gray}{#1}}
\newcommand{\sh}[1]{\cellcolor{gray!22}#1}
% Full-width strike line through the row just ended; place after \\.
\newcommand{\strikerow}{%
  \noalign{\vskip-0.5\dimexpr\ht\strutbox+\dp\strutbox\relax}%
  \noalign{\vskip-0.5\arrayrulewidth}%
  \arrayrulecolor{black!65}\hline\arrayrulecolor{black}%
  \noalign{\vskip 0.5\dimexpr\ht\strutbox+\dp\strutbox\relax}%
  \noalign{\vskip-0.5\arrayrulewidth}}
% Pointer to the long version (proofs and extended material).
\newcommand{\extended}{the full paper~\cite{hellerstein2026completecalmfull}}
\newcommand{\fullproof}{Proof in the full paper~\cite{hellerstein2026completecalmfull}.}
\newcommand{\proofinfull}{\begin{proof}In the full paper~\cite{hellerstein2026completecalmfull}.\end{proof}}

% ---------- Comment macro; NEVER EVER DELETE!!! ----------
\newcommand{\jmh}[1]{\textcolor{red}{\textit{#1 -- JMH}}}

\begin{document}

\title{Complete CALM: A Unified Criterion for Coordination}

\author{Joseph M. Hellerstein}
\affiliation{%
  \institution{UC Berkeley \& Amazon Web Services}
  \city{Berkeley, CA}
  \country{USA}
}

\begin{abstract}
When does a distributed system intrinsically need coordination?  The literature
has lacked a common formalism for the question.  CALM, CAP, and CRDTs answer it
for one class of programs, consistency models, or design patterns.  Their
answers cannot be compared or combined, and together they leave basic questions
open, such as how coordination repairs inconsistency.  We put them on common
ground.  We define a minimalist \spectype: a
possibility table that includes, for each input, all outputs a system may
produce.  Each column carries its own partial order, saying when one input
\inrel another and when one output \outrel another.  An \spectype is
\emph{consistent} when all its possible runs are confluent: runs whose inputs have a common
\indesc can always catch up to a common state.  Possibilities narrow as input
grows or outputs are emitted. An irrevocable \emph{\commitment} removes
possibilities not subsumed by those remaining; it requires coordination when
another run could contradict it.

In this unified setting, \emph{Functional CALM} generalizes
the CALM Theorem from Datalog queries to any function between ordered domains:
such a specification is consistent exactly when it is monotone: no
\commitment is required. We then go further: \emph{Complete CALM} handles nondeterministic specifications
that allow several outputs for the same input, where the test becomes
\emph{Lifted Monotonicity}: at every input, some permitted output \outrel
whatever could have been emitted up to that point.
\emph{Complete CAP} is a corollary: the CAP trade-off arises
exactly when Lifted Monotonicity fails.
\emph{Proper Coordination} identifies which stage of a composition must coordinate,
and certifies that later stages need no coordination of their own.  Highly Available Transactions, Invariant
Confluence, and later refinements of CALM all map into the same framework.
\end{abstract}

\begin{CCSXML}
<ccs2012>
<concept>
<concept_id>10003752.10003790.10003795</concept_id>
<concept_desc>Theory of computation~Distributed computing models</concept_desc>
<concept_significance>500</concept_significance>
</concept>
<concept>
<concept_id>10002951.10003317.10003347.10003356</concept_id>
<concept_desc>Information systems~Database transaction processing</concept_desc>
<concept_significance>300</concept_significance>
</concept>
</ccs2012>
\end{CCSXML}

\ccsdesc[500]{Theory of computation~Distributed computing models}
\ccsdesc[300]{Information systems~Database transaction processing}
\keywords{coordination, monotonicity, CALM, CAP, transactions, CRDTs}

\maketitle

% ======================================================================
\section{Introduction}
\label{sec:intro}

Concurrent systems rely on coordination mechanisms ranging from mutexes and
atomic instructions to database locks and consensus protocols. These mechanisms
prevent concurrent actions from producing undesirable outcomes.  They also bring
costs: overhead, latency, and availability in distributed settings.  Avoiding
unnecessary coordination is therefore a central goal in distributed systems and
databases.  But when is coordination avoidable, and when is it intrinsically
required? Various answers are scattered across the literature in incompatible
forms. We unify those results via a common formalism, and use it to address new
questions: when do nondeterministic specifications require coordination, what does
coordination achieve, and how do we know when it is properly applied?

The CALM Conjecture proposed a crisp, bidirectional criterion:
consistency can be achieved without coordination exactly when a specification
is monotone~\cite{hellerstein2010declarative}.  But the original proof applied
only in a restricted setting: Datalog queries evaluated by relational transducer
networks~\cite{ameloot2013relational,hellerstein2020keeping}. A related
question comes from the CAP conjecture: which consistency guarantees can remain
available when participants cannot communicate? There, the original proof
applied in a different restricted setting: linearizable
registers~\cite{brewer2000towards,gilbert2002cap}.  Both results are invoked
quite broadly: CALM as a general rule that monotone programs need no
coordination, CAP as a general constraint on distributed systems. The folklore
outgrew the proofs.

The apparent connection between CALM and CAP was never established in a common
formalism, and nearby results only complicate the picture. CRDTs use
inflationary updates over join-semilattices for coordination-free
convergence~\cite{shapiro2011crdt}; HATs identify transactional guarantees
available under partition~\cite{bailis2013hat}; invariant confluence asks when
independently reached database states merge
safely~\cite{bailis2014coordination}. The diverse formalisms leave cross-cutting
questions unclear, including whether CALM encompasses CAP, whether CRDT join
laws are necessary, why only some transactional isolation levels remain
available during partition, and whether invariant confluence differs from CALM
(\S\ref{sec:instances}).

\subsection{The Model: \Spectype[s]}
In systems, as in life,
possibilities narrow over time.  Sometimes the world moves on,
and possibilities are lost to us.  Other choices are left for us to decide, with
varying consequences.  Incidental decisions close no doors: each possibility they
remove is covered by some alternative. Fateful decisions
are the ones that matter: they exclude a possibility that nothing
remaining can cover.  Can we say which scenarios require
fateful decisions?

We ground the problem in a simple declarative model of \spectype[s],
independent of any language, protocol, or placement. The designer of a
specification provides a \emph{possibility relation} $P\subseteq
I\times O$: a table that may be infinite, like a mathematical function. Each row in
this table is a possibility observable from the outside: it pairs a cumulative input to a
permitted output. The design includes the input domain $I$ and output domain
$O$, with partial orders that help capture the specification's semantics.  On
inputs, $i\InputOrd j$ means that $j$ \emph{\inrel} $i$ with more input from the
environment.  On outputs, $u\OutputGe o$ means that $u$ subsumes all
specification-relevant information in $o$; we say
$u$ \emph{\outrel} $o$.  A \emph{\commitment} is an irrevocable, fateful system decision,
removing one or more rows $(i,o)$ for which no remaining row $(i,u)$ has
$o\OutputOrd u$: the output $o$ is no longer covered by any $u$,
and its possibility is gone.

State-machine-style specifications can be projected into this form, skipping
operational details that are not externally observable (\S\ref{sec:scope}).
Despite its simplicity, our model captures not only application invariants,
but also general-purpose consistency definitions such as linearizability,
causal consistency, and transactional isolation levels
(\S\ref{sec:complete-cap}, \S\ref{sec:hat-instance}, and \extended).

\subsection{The Criterion: Lifted Monotonicity}
With this model in hand, we are ready to define the key aspects of CALM:
Consistency and Monotonicity. A specification is \emph{consistent} when
all its possible executions are \emph{confluent}: states reached along
different paths can be brought to a common state. We will see that this definition
is general, accommodating narrower notions from prior literature.
Our monotonicity criterion requires more care: it must capture the full
expressivity of our model. It comes via two invariants. The first is always
required; the second addresses nondeterministic specifications.

\emph{1. Learn and refine.} As execution progresses, input grows: the system learns more from
the environment.
For specifications that permit only one output at each input, the possibility table is a function
and monotonicity is standard: $i\InputOrd j$ implies $f(i)\OutputOrd f(j)$.
\emph{Functional CALM} (\S\ref{sec:functional-calm-theorem}) proves that consistency is equivalent to
monotonicity in this setting.
% Emitting a monotone function's output is always an incidental choice:
% the output at any \inrelpast input is \outrelpart: it covers the earlier output.
% This setting generalizes Datalog-based CALM papers as well as CRDT state specifications.

% If monotonicity
% fails, some output requires a fateful decision: a \commitment that rules
% out later input which could contradict it.  The result covers the relational query
% languages of prior CALM papers, and the usual CRDT state specification.

\emph{2. Choose and align.} Prior work on CALM and CRDTs addresses functional
settings. \emph{Nondeterministic specifications}, including leader election and transaction
scheduling, are not functions: many outputs are within spec for a given input.
This requires an additional monotonicity invariant.  When many runs of a system
choose different outputs for the same input, those choices
must \outrelbase together---the outputs are \emph{directed}, in order theory terms.
Together, monotonicity under input growth and output choice make all output decisions incidental.
We call the pair of invariants \emph{Lifted Monotonicity}.
\emph{Complete CALM} proves that consistency is equivalent to Lifted
Monotonicity (\S\ref{sec:complete-calm-theorem}).

What of coordination?  A \commitment \emph{requires coordination}
if another run can contradict it; it is \emph{distributed coordination} if that
run is at another participant.
Distributed histories are captured under the same consistency criteria as centralized ones
(\S\ref{sec:distributed-execution}). Complete CAP (\S\ref{sec:complete-cap}) says
Lifted Monotonicity exactly characterizes specifications that can keep every
permitted output available and confluent during a partition.  The
Gilbert--Lynch theorem follows directly.

\paragraph{Contributions.}
We make contributions in three categories:
\begin{enumerate}
  \item \emph{A unified model.} The observable behavior of a system is captured via \spectype[s]: a possibility relation over
        ordered input and output domains. Consistency is defined as confluence, uniformly capturing general
        consistency models, application-specific requirements, and distributed consistency.
  \item \emph{A unified criterion.} For functions between ordered domains,
        Functional CALM shows that consistency is equivalent to monotonicity. For nondeterministic specifications,
        we prove Complete CALM: Consistency As Lifted Monotonicity. Complete CAP is a corollary:
        the tradeoff holds only under failure of Lifted Monotonicity.
  \item \emph{\Commitments, Composition and Proper Coordination.} We define \commitments to capture how
        systems exclude possibilities that are not covered by alternatives. We then show how specifications
        compose by equijoin, so that, in composition, upstream
        stages can free downstream stages from issuing \commitments.
\end{enumerate}

The full paper~\cite{hellerstein2026completecalmfull} gives further
intuition, proofs, and applications.

% ======================================================================
\section{Functional CALM}
\label{sec:functional-calm}

This section focuses on functional specifications, in which each input has a
single specified output.  We walk through the framework informally on an
election (\S\ref{sec:election-story}), state it precisely
(\S\ref{sec:framework-brief}), and
prove Functional CALM (\S\ref{sec:functional-calm-theorem}), identifying
when an output can be emitted without \commitment.

\subsection{The framework by example}
\label{sec:election-story}

\begin{figure*}[t]
  \centering
  \small
  \newcommand{\rowhdr}{\toprule $\mathrm{in}$ & $\mathrm{out}$ \\ \midrule}
  \setlength{\tabcolsep}{6pt}
  \begin{tabular}[t]{@{}l c@{}}
    \multicolumn{2}{@{}l}{(a) Initial state, current input $a\cdot\cdot$}\\
    \rowhdr
    \sh{$\mathbf{a\cdot\cdot}$} & \sh{$\bot$} \\
    $ab\cdot$ & \sh{$\bot$} \\
    $aba$ & \sh{$a$} \\
    $abb$ & \sh{$b$} \\
    $aa\cdot$ & \sh{$a$} \\
    $aaa$ & \sh{$a$} \\
    $aab$ & \sh{$a$} \\
    $a\cdot a$ & \sh{$a$} \\
    $a\cdot b$ & \sh{$\bot$} \\
    \bottomrule
    \multicolumn{2}{@{}l}{\footnotesize upper fiber $\{\bot,a,b\}$; undetermined}\\
    \multicolumn{2}{@{}l}{\footnotesize $\bot$ needs no \commitment}
  \end{tabular}
  \hfill
  \begin{tabular}[t]{@{}l c@{}}
    \multicolumn{2}{@{}l}{(b) After $\mathsf{input}(ab\cdot)$}\\
    \rowhdr
    \grayrow{$a\cdot\cdot$} & \grayrow{$\bot$} \\
    \sh{$\mathbf{ab\cdot}$} & \sh{$\bot$} \\
    $aba$ & \sh{$a$} \\
    $abb$ & \sh{$b$} \\
    \grayrow{$aa\cdot$} & \grayrow{$a$} \\
    \grayrow{$aaa$} & \grayrow{$a$} \\
    \grayrow{$aab$} & \grayrow{$a$} \\
    \grayrow{$a\cdot a$} & \grayrow{$a$} \\
    \grayrow{$a\cdot b$} & \grayrow{$\bot$} \\
    \bottomrule
    \multicolumn{2}{@{}l}{\footnotesize upper fiber $\{\bot,a,b\}$; undetermined}\\
    \multicolumn{2}{@{}l}{\footnotesize $\bot$ needs no \commitment}
  \end{tabular}
  \hfill
  \begin{tabular}[t]{@{}l c@{}}
    \multicolumn{2}{@{}l}{(c) After $\mathsf{input}(aa\cdot)$ instead}\\
    \rowhdr
    \grayrow{$a\cdot\cdot$} & \grayrow{$\bot$} \\
    \grayrow{$ab\cdot$} & \grayrow{$\bot$} \\
    \grayrow{$aba$} & \grayrow{$a$} \\
    \grayrow{$abb$} & \grayrow{$b$} \\
    \sh{$\mathbf{aa\cdot}$} & \sh{$a$} \\
    $aaa$ & \sh{$a$} \\
    $aab$ & \sh{$a$} \\
    \grayrow{$a\cdot a$} & \grayrow{$a$} \\
    \grayrow{$a\cdot b$} & \grayrow{$\bot$} \\
    \bottomrule
    \multicolumn{2}{@{}l}{\footnotesize upper fiber $\{a\}$; determined}\\
    \multicolumn{2}{@{}l}{\footnotesize $a$ needs no \commitment}
  \end{tabular}
  \hfill
  \begin{tabular}[t]{@{}l c@{}}
    \multicolumn{2}{@{}l}{(d) (b), then polls closed}\\
    \rowhdr
    \grayrow{$a\cdot\cdot$} & \grayrow{$\bot$} \\
    \sh{$\mathbf{ab\cdot}$} & \sh{$\bot$} \\
    \grayrow{$aba$} & \grayrow{$a$} \\ \strikerow
    \grayrow{$abb$} & \grayrow{$b$} \\ \strikerow
    \grayrow{$aa\cdot$} & \grayrow{$a$} \\
    \grayrow{$aaa$} & \grayrow{$a$} \\
    \grayrow{$aab$} & \grayrow{$a$} \\
    \grayrow{$a\cdot a$} & \grayrow{$a$} \\
    \grayrow{$a\cdot b$} & \grayrow{$\bot$} \\
    \bottomrule
    \multicolumn{2}{@{}l}{\footnotesize upper fiber $\{\bot\}$; determined}\\
    \multicolumn{2}{@{}l}{\footnotesize $\bot$ is the result}
  \end{tabular}
  \caption{The election specification and three states reached from it, showing rows
  whose input \inrel $a\cdot\cdot$.  Bold: current input; shaded: the current
  input and its upper fiber.  Gray: lost after an input event; struck
  through if removed by a \commitment.  Panels (b)
  and (c) follow input events.  In (d) the registrar emits $\bot$ and also
  applies a \commitment that $\mathsf{output}(\bot)$ alone does not require.}
  \label{fig:election-frames}
\end{figure*}

We introduce the framework with a small election, illustrated in
Figure~\ref{fig:election-frames}.  Each row of the specification table pairs
the ballots received so far with the permitted result; the column orders
say how inputs \ingrow and which later outputs \outrelbase earlier ones.

Three voters $v_1,v_2,v_3$ each cast at most one immutable ballot for $a$ or
$b$.  The output is whichever candidate has two ballots, or $\bot$ while there
is no winner.  An input is the ballots received so far, written as a
three-character string with a dot for a ballot not yet received.  For brevity,
the figure shows only inputs \inrelpart $a\cdot\cdot$.

The two orders are where a designer makes decisions.  The \emph{input order}
records accumulated ballots, one per voter: $ab\cdot$ \inrel $a\cdot\cdot$, while
$aa\cdot$ and $ab\cdot$ are incomparable.  The \emph{output order} says which
announcements are safe to build on: $\bot$ claims nothing, so either winner
\outrel it, while $a$ and $b$ are incomparable.  In general, the functional table is \emph{monotone}
when
\[
  i\InputOrd j\quad\Longrightarrow\quad f(i)\OutputOrd f(j).
\]
Here, for example, the input sequence $a\cdot\cdot\InputOrd ab\cdot\InputOrd aba$
monotonically produces the output sequence $\bot\OutputOrd\bot\OutputOrd a$.

\paragraph{Input and determination.}
Each panel captures the specification state: a pair of a \emph{current input}
and a \emph{remainder table} of rows
still possible from that input.  When $v_2$'s ballot for $b$ arrives, the input event
$\mathsf{input}(ab\cdot)$ produces panel~(b), causing the table to lose rows whose inputs no
longer \inrelbase the ballots received; such inputs are now illegal.  The \emph{upper fiber} is the set of
outputs in the remaining rows at the current input and its \indescs: in panel~(b),
$\{\bot,a,b\}$.
If $v_2$ instead votes for $a$ (panel~(c)), the upper fiber is $\{a\}$.  The state
is \emph{determined}: input alone resolved the outcome, without \commitment.

Panel~(d), by contrast, involves the \emph{system}: the implementation that
receives the ballots and announces the result
(\S\ref{sec:systems}).  In this example, the system acts as a registrar, emitting $\bot$
and closing the polls at
$ab\cdot$ in the same output event.  Its irrevocable \emph{\commitment}
removes the still-possible $aba$ and $abb$ rows but leaves the current input
unchanged, and so makes $\bot$ the final result.  Emitting $\bot$ by itself
would need no \commitment, since every output still possible \outrel it;
closing the polls is a further decision by the system, not something the
specification requires.

\begin{figure}[t]
  \centering
  \small
  \setlength{\tabcolsep}{6pt}
  \begin{minipage}{0.86\columnwidth}
  (e) As (b), all outputs incomparable, $\bot$ emitted\par\smallskip
  \begin{minipage}[t]{0.44\linewidth}
    \begin{tabular*}{\linewidth}[t]{@{}l@{\extracolsep{\fill}}c@{}}
      \toprule
      $\mathrm{in}$ & $\mathrm{out}$ \\ \midrule
      \grayrow{$a\cdot\cdot$} & \grayrow{$\bot$} \\
      \sh{$\mathbf{ab\cdot}$} & \sh{$\bot$} \\
      \grayrow{$aba$} & \grayrow{$a$} \\ \strikerow
      \grayrow{$abb$} & \grayrow{$b$} \\ \strikerow
      \grayrow{$aa\cdot$} & \grayrow{$a$} \\
    \end{tabular*}
  \end{minipage}\hfill
  \begin{minipage}[t]{0.44\linewidth}
    \begin{tabular*}{\linewidth}[t]{@{}l@{\extracolsep{\fill}}c@{}}
      \toprule
      $\mathrm{in}$ (cont.) & $\mathrm{out}$ (cont.) \\ \midrule
      \grayrow{$aaa$} & \grayrow{$a$} \\
      \grayrow{$aab$} & \grayrow{$a$} \\
      \grayrow{$a\cdot a$} & \grayrow{$a$} \\
      \grayrow{$a\cdot b$} & \grayrow{$\bot$} \\
      \bottomrule
    \end{tabular*}
    {\footnotesize upper fiber $\{\bot\}$; determined}
  \end{minipage}\par\smallskip
  \end{minipage}
  \caption{Panel (b) under the equality order, after $\mathsf{output}(\bot)$, wrapped
  to preserve space.  The
  The registrar's \commitment in Figure~\ref{fig:election-frames}(d) is required here.}
  \label{fig:nonmonotone-frame}
\end{figure}

\paragraph{Output and consistency.}
In a functional table, an output $o$ is a \emph{legal output} when every
output that could occur in future \outrel it---the entire upper fiber of the
current input.  Hence $\bot$ is legal in panels~(a) and~(b), and $a$ in panel~(c).  The atomic output event
$\mathsf{output}(o)$ empties any legal input whose output does not \outrelbase
$o$, then emits $o$, in one transition.  Emitting $o$ is a promise:
from then on, only rows whose outputs \outrelbase $o$ remain.  The emission
matters because the environment is powerful and unpredictable; it may use $o$
when choosing later input.

Figure~\ref{fig:nonmonotone-frame} considers a variant of our example that
changes only the output order: all outputs are incomparable, including $\bot$
for no winner (the \emph{equality order}, under which each element is related only
to itself).  Then $\bot$ is incompatible with future winners, so emitting it
requires a \commitment.  That change in the output order is Functional CALM in
miniature: the original order is monotone, so later outputs \outrelbase every
earlier one; under the equality order, some outputs can be emitted
only with a \commitment.

The obstruction is a property solely of the table and its orders.  Rows $(ab\cdot,\bot)$ and
$(aba,a)$ have the common input \indesc $aba$, but under the equality
order no output at $aba$ \outrel both outputs $\bot$ and $a$.  A run that emits
$\bot$ at $ab\cdot$ and one that receives $aba$ cannot be reconciled at $aba$.

Input and output events form a transition system, and the specification is
consistent when its runs are confluent.
Section~\ref{sec:functional-calm-theorem} shows that for functional tables
this corresponds exactly to monotonicity.

\subsection{The framework in brief}
\label{sec:framework-brief}
\label{sec:functional-specifications}
\label{sec:functional-events}
\label{sec:executions}

We now state these notions precisely.

\begin{definition}[Functional \spectype, total]
\label{def:functional-specification}
A \emph{functional \spectype} is a tuple
$\Spec=(I,\InputOrd,O,\OutputOrd,P)$,
where $(I,\InputOrd)$ and $(O,\OutputOrd)$ are partially ordered input and
output types, and $P\subseteq I\times O$ is its \emph{possibility relation}:
a typed relation with schema $(\mathrm{in},\mathrm{out})$ in which every
input appears in at least one row ($P$ is \emph{total} on $I$) and at most one
(a functional dependency).  The unique output paired with $i$ induces a
function $f:I\to O$.
\end{definition}

Because an output may depend on everything a system has received, we model
inputs as \emph{cumulative}: $i\InputOrd j$ means that $j$ contains at least
the input information in $i$.  We say that $j$ \emph{\inrel} $i$, or is an
\emph{\indesc} of $i$.  We call $\OutputOrd$ the \emph{output order}:
$o\OutputOrd u$ means that $u$ preserves every externally visible claim in $o$
and may add more, and we say that $u$ \emph{\outrel} $o$.  A \emph{common
\indesc} of inputs is an input that \inrel each of them, and a \emph{common
\outnoun} of outputs is an output that \outrel each of them.  Neither need be
unique or least: for example, two independent transactions $u$ and $v$, with schedules
ordered by subsequence, have two incomparable minimal common \indescs, $uv$
and $vu$.

\begin{definition}[Functional monotonicity, monotone]
\label{def:functional-monotonicity}
$\Spec$ is \emph{monotone} if
\begin{equation}
  i\InputOrd j
  \quad\Longrightarrow\quad
  f(i)\OutputOrd f(j).
  \label{eq:functional-monotonicity}
\end{equation}
\end{definition}

We introduce some compact notation. For a relation $R\subseteq I\times O$, a predicate in square brackets selects
rows, and commas separate conjuncts: $R[\mathrm{in}\InputGe j,\
\mathrm{out}\OutputGe o]$ holds the rows whose input \inrel $j$ and whose
output \outrel $o$.  We say that a row $(k,u)$ is \emph{at} $k$, and $\mathrm{out}(R)$ is
the set of outputs in $R$.

\paragraph{States.}
As illustrated by the non-gray text in the panels of Figure~\ref{fig:election-frames}, a \emph{specification state} is a pair $\sigma=(i,R)$ of a \emph{current
input} $i$ and a \emph{remainder table} $R\subseteq P$ that contains a row at
$i$ and only rows at \indescs of $i$.  An input $k$ is \emph{legal} at
$\sigma$ when $R$ contains a row at $k$.  Rows are never added during
execution.  The \emph{canonical state at $i$}, the state of a system that has
received $i$ and emitted nothing, has no rows removed: it keeps every row at $i$
and its \indescs:
\begin{equation}
  \sigma_P(i)=\bigl(i,\ P[\mathrm{in}\InputGe i]\bigr).
  \label{eq:canonical-state}
\end{equation}
The \emph{fiber} of $R$ at $k$ is $R(k)=\mathrm{out}(R[\mathrm{in}=k])$.  The
\emph{upper fiber} of $(i,R)$ is $\mathrm{out}(R)$, and the state is \emph{determined} when it has one element.

\paragraph{Events.}
For a state $(i,R)$ and a legal input $j$, the \emph{input event}
\begin{equation}
  (i,R)\xrightarrow{\mathsf{input}(j)}
  \bigl(j,\ R[\mathrm{in}\InputGe j]\bigr)
  \label{eq:input-event}
\end{equation}
makes $j$ current and \emph{subsumes} the rows whose input does not
\inrelbase $j$.  Because inputs are cumulative, $\mathsf{input}(j)$ then
$\mathsf{input}(k)$ has the effect of $\mathsf{input}(k)$ alone.

A system controls what it tells the world.  We study what it can
promise without assuming that clients tolerate retractions.  Later outputs
must therefore \outrelbase earlier ones.

\begin{definition}[Output event, canonical \commitment]
\label{def:output-event}
Let $(i,o)$ be the row at the current input of a state $(i,R)$.  The
\emph{output event} $\mathsf{output}(o)$ atomically restricts the remainder
table to the rows whose outputs \outrelbase $o$, and emits $o$:
\begin{equation}
  (i,R)\xrightarrow{\mathsf{output}(o)}\bigl(i,\ R[\mathrm{out}\OutputGe o]\bigr).
  \label{eq:output-action}
\end{equation}
When this removes a row, the transition applies the \emph{canonical \commitment}
$x_o$.  For every row it removes, no row remains at the same input.
\end{definition}

A \commitment removes rows that input alone would preserve; afterward, their
inputs are no longer legal.  Systems may issue more \commitments than an output
requires, as the registrar does in Figure~\ref{fig:election-frames}(d), but
extra \commitments only restrict behavior (\S\ref{sec:systems}).

\begin{definition}[Legal output]
\label{def:exposable}
The output $o$ at $(i,R)$ is a \emph{legal output} when $\mathsf{output}(o)$
applies no \commitment.
\end{definition}

In a functional specification, $o$ is a legal output exactly when $f(j)$
\outrel $o$ for every legal $j$.  An event from a determined state always
leads to another determined state.

\paragraph{Runs and consistency.}
We now consider the executions corresponding to a specification.
\begin{definition}[Specification event, run, \commitment-free run, reachable without \commitment]
\label{def:run}
Input and output events are \emph{specification events}, or just \emph{events}.  A \emph{run} is a
finite sequence of them from the
canonical state $\sigma_P(i)$ at some input $i$.  It is
\emph{\commitment-free} when no output event in it applies a \commitment; the
states such runs reach are \emph{reachable without \commitment}.
\end{definition}

Input events keep only rows whose input \inrel the current one, and output
events keep only rows whose outputs \outrelbase what has been emitted.  Both
steps are relational selections, and selections commute; since inputs only \ingrow during
a run, the last input's selection subsumes the earlier ones.

\begin{lemma}[Rows remaining after a run]
\label{lem:run-residual}
If a run ends at current input $i$ after emitting
$o_1,\ldots,o_m$, its remainder table is
$P[\mathrm{in}\InputGe i,\ \mathrm{out}\OutputGe o_1,\ \ldots,\
  \mathrm{out}\OutputGe o_m]$.
\end{lemma}
\proofinfull

To compare runs that end at different inputs, we consider a common \indesc
of those inputs.

\begin{definition}[Consistency, confluence at an input]
\label{def:consistency}
An \spectype is \emph{consistent} when every finite family of
runs is confluent: if the runs end at states
$(i_1,R_1),\ldots,(i_k,R_k)$ and $j$ is a common \indesc of their current
inputs, each can be continued to a
common state with current input $j$.  They are then \emph{confluent at} $j$.
\end{definition}

Consistency requires confluence at every common \indesc, not merely at some
later one.  The reason is observability: the environment may act on conflicting
outputs through channels the system does not see, as in Lamport's ``anomalous
behavior''~\cite{lamport1978time}, and a reconciling input may never arrive.
Thus a conflict is a consistency failure even if later input could resolve it.

Confluence has a simple test: some row at the common \indesc must survive
in every run.

\begin{proposition}[Survivor test]
\label{prop:confluence}
Runs ending at states $(i_r,R_r)$, $r=1,\ldots,k$, are confluent
at $j$ if and only if some row $(j,u)$ at $j$ is in every $R_r$.  For any such
$u$, $\mathsf{input}(j)$ followed by $\mathsf{output}(u)$ takes every final
state to $\bigl(j,\ P[\mathrm{in}\InputGe j,\ \mathrm{out}\OutputGe u]\bigr)$.
\end{proposition}

\begin{proof}[Proof sketch]
Rows are never added, so the row of a common state at $j$ was there
in every run.  \fullproof
\end{proof}

Consistency guarantees runs never diverge irreparably: no ``split brain''.  Whatever input
and output events a finite collection of runs has taken, any input \inrelpart all
of theirs leaves a single state that every run can reach: one that keeps every
promise any of them has made and accepts all of that input.
Consistency is a property of the specification.  It describes no protocol and
promises no delivery or termination.

\subsection{The Functional CALM theorem}
\label{sec:functional-calm-theorem}

Consistency is defined by a specification's executions, monotonicity by its
table and orders alone.  Functional CALM shows that these coincide.  Full proofs are in \extended.

\begin{theorem}[Functional CALM]
\label{thm:functional-calm}
A functional \spectype is consistent if and only if it is
monotone.
\end{theorem}

\begin{proof}[Proof sketch]
A monotone $f$ only emits outputs that all of its future outputs \outrelbase,
so the row $(j,f(j))$ at a common \indesc $j$ survives in every run.
If $f$ is not monotone, a run that emits $f(i)$ and a run that accepts an
\indesc $k$ of $i$ with $f(i)\not\OutputOrd f(k)$ can never be reconciled at $k$.
\fullproof
\end{proof}

In the equality-order election, $f(ab\cdot)=\bot$ and $f(aba)=a$ are
incomparable: a run that emits $\bot$ at $ab\cdot$ and one that accepts
$aba$ are the two runs of Section~\ref{sec:election-story} that can never be reconciled.

\begin{corollary}[State is fixed by input]
\label{cor:state-by-input}
In a consistent functional specification, every run ending at current
input $i$ is in the canonical state $\sigma_P(i)$, whatever outputs it emitted
along the way.
\end{corollary}
\proofinfull

Knowing the input is therefore enough to know the state, and with it the one
output.  Monotonicity also governs legality:

\begin{proposition}[Output-legality boundary]
\label{prop:functional-exposure-boundary}
A functional \spectype is monotone if and only if every output at
every state reachable without \commitment is legal.
\end{proposition}

\begin{proof}[Proof sketch]
By Corollary~\ref{cor:state-by-input}, monotone $f$ reaches only canonical
states. At $\sigma_P(i)$, $f(i)$ is legal exactly
when $f$ is monotone at $i$.
 \fullproof
\end{proof}

A monotone specification can thus emit every current
output without \commitment; if monotonicity fails at $i\InputOrd k$ with
$f(i)\not\OutputOrd f(k)$, then at $\sigma_P(i)$ the event $\mathsf{output}(f(i))$
applies $x_{f(i)}$, which excludes input $k$.

This is the familiar CALM dichotomy stated at the level of the specification
rather than inside a particular programming language.  The C in CALM has been read
both as consistency and as confluence; Definition~\ref{def:consistency} equates them.
Classical CALM fits the pattern with inputs $J$ as sets of facts ordered by
inclusion and output tables $Q(J)$ ordered by inclusion.  Functional
monotonicity becomes $J\subseteq J'\Rightarrow Q(J)\subseteq Q(J')$,
the standard definition of a monotone query~\cite{abiteboul1995foundations},
which positive Datalog satisfies and negation may not.  Functional CALM
generalizes the statement to arbitrary input and
output orders, whatever formalism expresses the specification.  CRDT state specifications have the same functional shape,
with received payloads ordered by inclusion and states by the semilattice
order~\cite{shapiro2011crdt}; Sections~\ref{sec:query-instance}
and~\ref{sec:crdt-instance} treat both.

% ======================================================================
\section{Complete CALM}
\label{sec:complete-calm}

We now extend CALM theory to allow a fiber to contain several outputs.  The monotonicity property we want to test as input
\ingrows from $i$ to $j\InputGe i$ must be lifted to sets: can \emph{some}
output in the fiber at $j$ \outrelbase \emph{whatever} outputs runs emit at $i$?
A \commitment can empty an input's fiber or remove some output rows while it
remains nonempty.  We now turn to the latter effect.  Section~\ref{sec:ambiguity-story}
introduces the issue by example.  Section~\ref{sec:lifted-monotonicity}
defines \spectype[s] and Lifted Monotonicity.  Section~\ref{sec:lifted-output-events} lifts
output events and classifies what they remove, and
Section~\ref{sec:complete-calm-theorem} shows that the run-level results of
Section~\ref{sec:executions} carry over and proves Complete CALM.

\subsection{\Commitment by example}
\label{sec:ambiguity-story}

Figure~\ref{fig:ambiguity-frames} contrasts two specifications that permit
several outputs in one fiber.  Transaction ordering requires a \commitment
between incompatible serialization orders.  For the discussion here, outputs are simply serial
schedules, to which executions correspond via equivalence; Section~\ref{sec:hat-instance} studies more thorough
transactional specifications.  Traceroute observations, by
contrast, can all be \outrelpast by a larger route set.

\begin{figure}[t]
  \centering
  \small
  \setlength{\tabcolsep}{4pt}
  \begin{minipage}[t]{0.36\columnwidth}
    \raggedright
    (f) After $[A,B]$\par\smallskip
    \begin{tabular*}{\linewidth}[t]{@{}l@{\extracolsep{\fill}}l@{}}
      \toprule
      $\mathrm{in}$ & $\mathrm{out}$ \\ \midrule
      $\{A,B\}$ & \sh{$[A,B]$} \\
      \grayrow{$\{A,B\}$} & \grayrow{$[B,A]$} \\ \strikerow
      \bottomrule
    \end{tabular*}\par\smallskip
    {\footnotesize no remaining row \outrel the removed one: a \commitment}
  \end{minipage}\hfill
  \begin{minipage}[t]{0.6\columnwidth}
    \raggedright
    (g) After $\{s{\to}\texttt{s1}\}$\par\smallskip
    \begin{tabular*}{\linewidth}[t]{@{}l@{\extracolsep{\fill}}l@{}}
      \toprule
      $\mathrm{in}$ & $\mathrm{out}$ \\ \midrule
      $i$ & \sh{$\{s{\to}\texttt{s1}\}$} \\
      \grayrow{$i$} & \grayrow{$\{s{\to}h\}$} \\
      \grayrow{$i$} & \grayrow{$\{s{\to}h,\ s{\to}h{\to}\texttt{s2}\}$} \\
      $i$ & \sh{$\{s{\to}\texttt{s1},\ s{\to}h\}$} \\
      $i$ & \sh{$\{s{\to}\texttt{s1},\ s{\to}h,\ s{\to}h{\to}\texttt{s2}\}$} \\
      \bottomrule
    \end{tabular*}\par\smallskip
    {\footnotesize each removed row is \outrelpast by a remaining one: an output
    subsumption}
  \end{minipage}
  \caption{Two specifications with several rows at one input, each after one output
  event.  (f) Serial orders of two concurrent transactions, under the equality
  order.  (g) Route sets from a source $s$ to a service on \texttt{s1} and
  \texttt{s2}, ordered by inclusion; $s$ reaches \texttt{s1} directly and
  \texttt{s2} through hop $h$.  Legend as in Figure~\ref{fig:election-frames};
  here gray alone marks output subsumption.}
  \label{fig:ambiguity-frames}
\end{figure}

\paragraph{\Commitment and directedness.} The examples differ in what happens
to alternative outputs when one is emitted.  In panel~(f), neither
serialization order \outrel the other, so emitting one requires an
\emph{\commitment} that removes the other while the input stays legal:
emitting $[A,B]$ excludes $[B,A]$.  In panel~(g), larger route sets \outrelbase both the
emitted observation and the outputs of the rows it removes.  We call such
removal an \emph{output subsumption}; every removed output is
\outrelpast by a remaining one, so emitting the route observation loses no information
according to the output order and needs no
\commitment: it is a \emph{legal output}.

An output event can thus remove rows in two ways: by \commitment, as in panel~(f) and Section~\ref{sec:functional-calm}, or by subsuming outputs, as in panel~(g).  At a given input, a \commitment may remove every row, ruling out the input, or only some rows, ruling out those outputs.  The atomic event $\mathsf{output}(o)$ does whichever $o$ requires, then emits $o$.  Subsuming is not \commitment.

In Figure~\ref{fig:ambiguity-frames}, serialization requires \commitment; traceroute avoids it.
The difference between them is the order-theoretic concept of a \emph{directed} fiber:
every finite set of outputs in a given fiber has a common \outnoun in that fiber.
Directedness at a single input is not sufficient, however;
we must also \outrelbase prior outputs as input \ingrows. \emph{Lifted Monotonicity}
combines the requirements: as input \ingrows from $i$ to $j$, some output
in the fiber at $j$ \outrel whatever outputs runs emit at $i$.  Equivalently,
outputs emitted at inputs with a common \indesc $j$ must have a common
\outnoun in the fiber at $j$: the same test as in Section~\ref{sec:election-story}, but
with many outputs per input.  Two transaction scheduling runs that emit $[A,B]$ and $[B,A]$ fail it;
two traceroute runs that emit different route sets pass it, since their union
survives in both.

\subsection{Lifted Monotonicity}
\label{sec:relational-specifications}
\label{sec:lifted-monotonicity}

We begin with our most general definition of a specification, which admits
output nondeterminism.
\begin{definition}[\Spectype, nondeterminism]
\label{def:relational-specification}
An \emph{\spectype} is
$\Spec=(I,\InputOrd,O,\OutputOrd,P)$ with $P\subseteq I\times O$,
where $(I,\InputOrd)$ and $(O,\OutputOrd)$ are partially ordered sets.  The
relation $P$ is the specification's \emph{possibility relation} and is total
on $I$.  A row $(i,o)\in P$ is a possible input/output pair: output $o$ is
permitted at input $i$.  The fiber
$\fiber{i}=\{o\in O:(i,o)\in P\}$ may now contain several
outputs; the \spectype is \emph{nondeterministic} when some fiber does.  A
functional \spectype is the special case $\fiber{i}=\{f(i)\}$.  Specification
nondeterminism describes permitted results, not a strategy's method of choosing
among them, and does not by itself imply a \commitment.
\end{definition}

Specification states and input events
(\S\ref{sec:framework-brief})
apply unchanged to relational tables; Section~\ref{sec:lifted-output-events} lifts
output events.

As input
\ingrows from $i$ to $j$, can some output in the fiber at $j$ \outrelbase
whatever outputs runs emit at $i$?  Two things can go wrong.  A
\emph{stranded output} is an output in the fiber at $i$ with no \outnoun in
the fiber at some \indesc $j$.  An \emph{incompatible pair} is two outputs in
the fiber at $i$ with no common \outnoun there.  We rule out
each, combine them, and restate the result in the form runs need.

First, we rule out stranded outputs.

\begin{definition}[Relational Monotonicity]
\label{def:relational-monotonicity}
An \spectype satisfies \emph{Relational Monotonicity}
when every output keeps some \outnoun as input \ingrows:
        \begin{equation}
  i\InputOrd j,\ o\in\fiber{i}
  \quad\Longrightarrow\quad
  \exists u\in\fiber{j}:\ o\OutputOrd u.
  \label{eq:relational-monotonicity}
        \end{equation}
\end{definition}

For singleton fibers this is functional monotonicity,
Equation~\ref{eq:functional-monotonicity}.  It treats each output separately,
so it does not rule out incompatible pairs, which concern several outputs in
one fiber.
For a set $F\subseteq O$ of outputs, outputs $o_1,\ldots,o_k$ are
\emph{compatible in $F$} if some $u\in F$ \outrel all of them, and
\emph{incompatible in $F$} otherwise.

Second, we rule out incompatible pairs.

\begin{definition}[Directed fiber, Directed Outputs]
\label{def:fiber-directedness}
A fiber $F\subseteq O$ is \emph{directed} when every nonempty finite
$X\subseteq F$ is compatible in $F$.  An \spectype has
\emph{Directed Outputs} when every fiber $\fiber{i}$ is directed.
\end{definition}

A singleton fiber is trivially directed.  As with common \indescs of inputs, a
common \outnoun need not be
least: directedness requires some upper bound within the fiber for each finite
set, not a least one or a fixed merge operator, so it is weaker than a
join-semilattice.  For example, the fiber
\[
  \{\{a\},\{b\},\{a,b,c\},\{a,b,d\},\{a,b,c,d\}\}
\]
under inclusion is directed, since $\{a,b,c,d\}$ \outrel every member, yet $\{a\}$ and $\{b\}$ have two
incomparable minimal common \outnouns.  Thus
the traceroute fiber is directed: the union of finitely many route sets in
it is again in the fiber and refines each. The two transaction orders in
Figure~\ref{fig:ambiguity-frames}(f), by contrast, are incompatible.

\paragraph{Lifted Monotonicity.}
To state both requirements at once, lift $\OutputOrd$ to sets.  For nonempty
$A,B\subseteq O$, define
        \begin{equation}
  A\LiftedOrd B
  \;\Longleftrightarrow\;
  \text{every nonempty finite }X\subseteq A\text{ is compatible in }B,
  \label{eq:joint-output-order}
        \end{equation}
the \emph{joint lift} of $\OutputOrd$.  We apply it with $A=\fiber{i}$ and
$B=\fiber{j}$ for $i\InputOrd j$.  When $X$ is a single output $o$, it asks
that $o$ have a \outnoun in $\fiber{j}$, which rules out stranded outputs.  When $X$
holds several outputs, it asks that they have a common \outnoun in $\fiber{j}$; with
$j=i$, this rules out incompatible pairs.

\begin{definition}[Lifted Monotonicity, monotone specification]
\label{def:lifted-monotonicity}
An \spectype satisfies \emph{Lifted Monotonicity}, and
is called \emph{monotone}, when the mapping $i\mapsto\fiber{i}$ is monotone
under the joint lift:
        \begin{equation}
  i\InputOrd j
  \quad\Longrightarrow\quad
  \fiber{i}\LiftedOrd\fiber{j}.
  \label{eq:lifted-monotonicity}
        \end{equation}
\end{definition}

With singleton fibers it is functional monotonicity
(Definition~\ref{def:functional-monotonicity}). Lifted
Monotonicity holds exactly when Relational Monotonicity and Directed Outputs
both do.

\paragraph{The common-\indesc condition.}
Lifted Monotonicity compares an input with its
own expansions. Runs may emit outputs at incomparable inputs, say
$o_1$ at $i_1$ and $o_2$ at $i_2$, and can be reconciled only at a common
\indesc $j$ of both. The next condition asks exactly that: a common
\outnoun of $o_1$ and $o_2$ in the fiber at $j$.

\begin{definition}[Common-\indesc condition]
\label{def:common-expansion-condition}
An \spectype satisfies the \emph{common-\indesc condition} when,
for every finite set of rows $(i_1,o_1),\ldots,(i_k,o_k)$ of $P$ and every
common \indesc $j$ of $i_1,\ldots,i_k$, some $u\in P(j)$ \outrel every
$o_r$.
\end{definition}

For the empty set of rows the condition asks only that $P(j)$ be nonempty, which
follows from $P$ being total on $I$.

\begin{proposition}[Common-\indesc form of Lifted Monotonicity]
\label{prop:common-expansion-form}
An \spectype satisfies Lifted Monotonicity if and only
if it satisfies the common-\indesc condition.
\end{proposition}

\begin{proof}[Proof sketch]
Relational Monotonicity carries each $o_r$ to some
\outnoun in the fiber at $j$, and Directed Outputs gives those a common
\outnoun there.  \fullproof
\end{proof}

\subsection{Lifted output events}
\label{sec:lifted-output-events}

Lifted Monotonicity is a property of the possibility relation and its orders.  We now return to execution.
For a state $(i,R)$ and $o\in R(i)$, the output event $\mathsf{output}(o)$
again keeps exactly the rows whose outputs \outrelbase $o$, ending at
$(i,R[\mathrm{out}\OutputGe o])$ as in Equation~\ref{eq:output-action}, and
emits $o$.  Its removed rows fall into two classes.  A row is in
$\Subsumed_o$ when a kept row at the same input \outrel it; its removal is an
\emph{output subsumption}.  Every other removed row is in $\Excluded_o$ and is
removed by the canonical \commitment $x_o$.  At an input $k$, $x_o$ rules out $k$ if it removes every row there; otherwise
it rules out each output whose row it removes at $k$.  Output subsumption is not \commitment.
Emitting $\bot$ under the equality order of Figure~\ref{fig:nonmonotone-frame}
puts the rows for $aba$ and $abb$ in $\Excluded_o$ and rules out those inputs;
emitting $[A,B]$ in Figure~\ref{fig:ambiguity-frames}(f) puts the row for
$[B,A]$ in $\Excluded_o$ and rules out that output at $\{A,B\}$; and the
route sets removed in panel~(g) are in $\Subsumed_o$.  In the functional
case, $\Subsumed_o$ is empty.  Formal definitions are in \extended.

As before, $o$ is a legal output when $\mathsf{output}(o)$ applies no
\commitment, exactly when $\Excluded_o$ is empty.
Equivalently, at every legal \indesc $j$ of $i$, $o$ and each
remaining output at $j$ have a common \outnoun at $j$.

\subsection{The Complete CALM theorem}
\label{sec:complete-calm-theorem}

A lifted output event ends at $(i,R[\mathrm{out}\OutputGe o])$, as in
Equation~\ref{eq:output-action}.  States, input events, runs,
\commitment-free runs, and consistency are defined as before, and the results
of Section~\ref{sec:executions} used output events only through that
equation.  They lift unchanged: the rows remaining after a run
(Lemma~\ref{lem:run-residual}) and the survivor test
(Proposition~\ref{prop:confluence}).  Thus consistency
reduces to Lifted Monotonicity.

\begin{theorem}[Complete CALM]
\label{thm:complete-calm}
An \spectype is consistent if and only if it satisfies Lifted
Monotonicity, equivalently the common-\indesc condition.
\end{theorem}

\begin{proof}[Proof sketch]
If the common-\indesc condition holds, a common \outnoun $u\in P(j)$
of everything emitted survives in every run.
If outputs emitted at inputs with a common \indesc $j$ are incompatible in
$P(j)$, the runs that emit them share no row at $j$ and can never be
reconciled there.  \fullproof
\end{proof}

The name CALM, for Consistency As Logical Monotonicity, tied the idea to a
logic language.  The new, more general expansion is Consistency As Lifted Monotonicity.
The condition that characterizes consistency is now a
property of a specification: Lifted Monotonicity, which for a function is
ordinary monotonicity.  Section~\ref{sec:distributed-consistency} shows that
both theorems hold for distributed histories.

The same condition governs legality.

\begin{proposition}[Output-legality boundary]
\label{prop:exposure-boundary}
An \spectype satisfies Lifted Monotonicity if and only
if every output at every state reachable without \commitment is legal.
\end{proposition}

\begin{proof}[Proof sketch]
Under the common-\indesc condition, consistency supplies the common \outnoun
the legal-output test asks for.  Conversely, the legal-output test at canonical
states, one row at a time, builds the common \outnoun the condition
asks for.  \fullproof
\end{proof}

For functions, both are monotonicity, so Complete CALM and
Proposition~\ref{prop:exposure-boundary} reduce to Functional CALM and
Proposition~\ref{prop:functional-exposure-boundary}.

% ======================================================================
\section{Distributed Execution}
\label{sec:distributed-execution}

Consistency compares independent runs.  A distributed execution places one run
at each participant, and the only change from our framework's perspective is
that participants can learn about inputs at other participants via messages.
Functional CALM and Complete CALM therefore have a natural distributed reading,
which Section~\ref{sec:histories} makes precise.  We then define systems, whose
\commitments may exceed those the specification requires, and \commitments that require distributed coordination, which a run at
another participant could contradict; Complete CAP follows.

\subsection{Distributed histories and consistency}
\label{sec:histories}
\label{sec:distributed-consistency}

A finite set $P$ of \emph{participants} each holds a specification state,
starting at a canonical state $\sigma_P(i_p)$, and changes it only by its own
specification events.  A participant learns input directly from the
environment or in a message from another participant.  A
\emph{message} $m=(j,d)$ carries a cumulative input $j$, or $\varnothing$
if it carries none, and arbitrary application-specific data $d$; one can envision timestamps, version
vectors, state identities, and so on.  A receive does not change the specification state; a
participant at $(i,R)$ that receives $(j,d)$ may incorporate $j$ by its own
$\mathsf{input}(k)$ for a legal common \indesc $k$ of $i$ and $j$.  Data can
influence a system's subsequent actions (\S\ref{sec:systems}), but never directly changes the state.

\begin{definition}[Distributed history, history event, happens-before, consistent cut]
\label{def:distributed-history}
A \emph{distributed history} $H$ is a finite set of
\emph{history events}, each an occurrence at one participant of a
specification event, $\mathsf{send}(m)$, or $\mathsf{receive}(m)$, with
a total local order on each participant's events, a matching send at another
participant for each receive, and, for each participant $p$, specification
events that in local order form a defined transition sequence from
$\sigma_P(i_p)$.  A \emph{consistent cut} $C$ is a
subset of history events closed under
happens-before~\cite{lamport1978time,chandy1985distributed} (a term
unrelated to our consistency).
\end{definition}

At a consistent cut $C$, the specification events of $C$ at participant $p$,
in local order, form its \emph{participant run}.  The cut is
\emph{confluent at} $j$, for a common \indesc $j$ of the participants'
current inputs, when some continuation of $C$ brings every participant to the
same state $(j,R)$.  A specification has \emph{distributed consistency} when,
for every finite participant set, every consistent cut of every history is
confluent at every such $j$.

\begin{proposition}[Consistency and distributed consistency]
\label{prop:distributed-consistency}
Consistency is equivalent to distributed consistency.
\end{proposition}

\begin{proof}[Proof sketch]
The participant runs at a cut are a family of runs, and any family
of runs can be placed at separate participants that exchange no messages.  \fullproof
\end{proof}

By Complete CALM, then, a specification has distributed consistency exactly
when it satisfies Lifted Monotonicity.  Under Lifted Monotonicity, this
confluence has a concrete witness.  For participant runs with common \indesc $j$,
the survivor test supplies a $u\in\fiber{j}$ that survives in every endpoint;
by Lemma~\ref{lem:run-residual}, $u$ \outrel everything they have emitted.  At
each participant, $\mathsf{input}(j)$ followed by $\mathsf{output}(u)$ reaches
the same state without \commitment (Proposition~\ref{prop:exposure-boundary}).  This is an existence witness in the
transition system; choosing $u$, communicating it, and scheduling the
transitions belong to the operational model (\S\ref{sec:scope}).

For functional specifications,
Corollary~\ref{cor:state-by-input} also gives convergence: participants with
the same current input are in the same state, regardless of message paths,
arrival order, or the outputs they emitted.  This is the convergence that strong eventual consistency~\cite{shapiro2011crdt}
asks of replicated state.  The semilattice merge of CRDTs suffices for it but
is not needed: here it follows directly from any consistent functional
specification.
With several outputs per input this fails even for consistent specifications:
in Figure~\ref{fig:ambiguity-frames}(g), a run that emitted
$\{s{\to}\texttt{s1}\}$ has subsumed the route sets that omit it, while a run at
the same input that emitted nothing keeps them.  Consistency then guarantees
that such replicas can be reconciled, not that they agree.

\subsection{Systems}
\label{sec:systems}

\begin{figure}[t]
  \centering
  \small
  \setlength{\tabcolsep}{5pt}
  \begin{tabular}[t]{@{}l l@{}}
    \multicolumn{2}{@{}l}{(h) 2PL: $A$ locks ahead of $B$}\\
    \toprule
    $\mathrm{in}$ & $\mathrm{out}$ \\ \midrule
    $\{A,B,C\}$ & \sh{$[A,B,C]$} \\
    $\{A,B,C\}$ & \sh{$[A,C,B]$} \\
    $\{A,B,C\}$ & \sh{$[C,A,B]$} \\
    \grayrow{$\{A,B,C\}$} & \grayrow{$[B,A,C]$} \\ \strikerow
    \grayrow{$\{A,B,C\}$} & \grayrow{$[B,C,A]$} \\ \strikerow
    \grayrow{$\{A,B,C\}$} & \grayrow{$[C,B,A]$} \\ \strikerow
    \bottomrule
  \end{tabular}
  \hfill
  \begin{tabular}[t]{@{}l l@{}}
    \multicolumn{2}{@{}l}{(i) MVTO: $B$ receives the first timestamp}\\
    \toprule
    $\mathrm{in}$ & $\mathrm{out}$ \\ \midrule
    $\{A,B,C\}$ & \sh{$[B,A,C]$} \\
    $\{A,B,C\}$ & \sh{$[B,C,A]$} \\
    \grayrow{$\{A,B,C\}$} & \grayrow{$[A,B,C]$} \\ \strikerow
    \grayrow{$\{A,B,C\}$} & \grayrow{$[A,C,B]$} \\ \strikerow
    \grayrow{$\{A,B,C\}$} & \grayrow{$[C,A,B]$} \\ \strikerow
    \grayrow{$\{A,B,C\}$} & \grayrow{$[C,B,A]$} \\ \strikerow
    \bottomrule
  \end{tabular}
  \caption{Serial orders of three transactions after one \commitment:
  (h) 2PL orders $A$ before $B$; (i) MVTO places $B$ first.  Only complete rows without aborts or reads
  are shown.}
  \label{fig:cc-commitments}
\end{figure}

Some systems apply \commitments before an output needs them.  The registrar can
emit $\bot$ without closing the polls, while a lock manager can fix an order
between ongoing transactions without emitting it.  We therefore let systems
bring bespoke \commitments.  In general, a \emph{\commitment} is a fateful
decision: a partial function on specification states that removes a nonempty
set of rows.  No row it removes is subsumed by a row that remains at the same
input, and it leaves the current input unchanged.  At an input $k$, a
\commitment may remove every row, ruling out $k$, or only some rows, ruling out
those outputs.  An output event
$\Out{\gamma}{o}$ applies a sequence $\gamma$ of \commitments, subsumes rows,
and emits $o$.
A \emph{\commitment event} $\mathsf{apply}(\gamma)$ applies such a sequence and
emits nothing, as when 2PL orders $A$ before $B$ while $B$ waits for $A$'s
lock; it is a third kind of specification event.  A \emph{system} consists of
a \emph{\commitment basis} $\mathcal B$ (a set of \commitments) and a
nondeterministic \emph{strategy} per participant that picks its next output,
\commitment, send, or receive event from its local history, applying only
\commitments from $\mathcal B$.  The full paper gives these definitions formally and shows
that every canonical \commitment of Section~\ref{sec:complete-calm} is a \commitment
in this sense.

Two-phase locking (2PL) and multiversion timestamp
ordering (MVTO)~\cite{bernstein1981concurrency} have different bases: 2PL's
$\mathsf{order}(A,B)$ fixes $A$ before $B$ and nothing else, while MVTO's
$\mathsf{stamp}(B)$ fixes $B$ before any transaction not yet stamped
(Figure~\ref{fig:cc-commitments}).  Both are \commitments, but they leave
different alternatives open.  We call the runs of
Sections~\ref{sec:functional-calm}--\ref{sec:complete-calm} \emph{canonical
runs}.  Locks, timestamps, quorums, and consensus differ in how and when
they decide; what the environment sees does not, as the next lemma shows.

\begin{lemma}[Canonical Run Universality]
\label{lem:canonical-equivalent}
Every run under any \commitment basis is \emph{observably equivalent} to a
canonical run: one that accepts the same inputs, emits the same outputs,
and keeps at least the original's rows after each event (\extended).
\end{lemma}

The canonical event for a legal output applies no \commitment, and deferring
that output avoids none.  To
illustrate, when 2PL orders $A$ before $B$ in a \commitment event, the
canonical run skips it.  If 2PL later emits $[A,C,B]$, the canonical run can
emit the same output because that row was present all along.  Every surviving
refinement of $[A,C,B]$ already orders $A$ before $B$, so the earlier \commitment
makes no further difference to the final table.  Whether a specification
needs coordination thus depends on the \commitments its outputs
require, whatever mechanism makes them.

\subsection{Coordination}
\label{sec:coordination}

Having separated a system's observable behavior from when and how it makes
\commitments, we now ask which \commitments require coordination.  The idea is
simple: a \commitment requires coordination if another run can contradict it.

\begin{definition}[Contradicting run]
\label{def:contradicts}
Let an event of a run $r$ apply a \commitment.  A run $r'$ \emph{contradicts}
that \commitment at input $k$ if it accepts $k$ after the \commitment emptied
its fiber, or emits at $k$ an output whose row the \commitment removed.
\end{definition}

We distinguish the behavior permitted by a specification from the behavior
admitted by a particular system.  Write $\mathcal R_P^\top$ for the set of all
canonical runs of $P$.  By Lemma~\ref{lem:canonical-equivalent}, it is maximal
at the observable interface: every run under any \commitment basis has the
same input and output events as some run in $\mathcal R_P^\top$.  A system $A$
implementing $P$ may admit fewer such behaviors because its strategies
restrict which events it can perform; write $\mathcal R_A$ for the runs it
produces.

\begin{definition}[Coordination-requiring \commitment, distributed coordination]
\label{def:coordination}
For any set of runs $\mathcal R$, a \commitment \emph{requires coordination}
relative to $\mathcal R$ when some run in $\mathcal R$ contradicts it
(Definition~\ref{def:contradicts}).  If participant $p$ applies the \commitment
in a history, it \emph{requires distributed coordination} when a run at
another participant $q\ne p$ of that history contradicts it.
\end{definition}

\begin{definition}[Coordination-free]
\label{def:coordination-free}
A run or history is \emph{coordination-free} relative to $\mathcal R$ when it
applies no \commitment that requires coordination relative to $\mathcal R$.
\end{definition}

\begin{proposition}[Canonical runs contradict every \commitment]
\label{prop:canonical-contradiction}
Every \commitment applied in a run from a canonical state $\sigma_P(i)$ is
contradicted by a canonical run from $\sigma_P(i)$.
\end{proposition}

\begin{proof}[Proof sketch]
Choose a row $(k,x)$ that the \commitment removes.  If it empties the fiber at
$k$, the canonical run that accepts $k$ contradicts it.  Otherwise, the
canonical run that accepts $k$ and emits $x$ contradicts it.
\end{proof}

\begin{corollary}[Canonical coordination]
\label{cor:canonical-coordination}
Relative to $\mathcal R_P^\top$, a run is coordination-free if and only if it
applies no \commitment.
\end{corollary}

\begin{proof}
Every applied \commitment has a contradicting run in $\mathcal R_P^\top$ by
Proposition~\ref{prop:canonical-contradiction}; the converse is immediate.
\end{proof}

Output subsumption remains permitted.  Given two distinct participants both
initialized at $i$, the run applying a \commitment from $\sigma_P(i)$ and a
canonical run that contradicts it can occupy separate participants in a
message-free history, by the construction in the proof of
Proposition~\ref{prop:distributed-consistency}.  The \commitment therefore
requires distributed coordination relative to that history's runs.

\begin{proposition}[Contradicting runs cannot reconcile at the conflict input]
\label{thm:commitments-require-coordination}
Consider a run that applies a \commitment and another run that contradicts it
at input $k$.  Their resulting states have no common continuation to a state
$(k,R')$.
\end{proposition}

\begin{proof}[Proof sketch]
If the contradicting run accepts $k$, the \commitment emptied its fiber; since
events never add rows, the first run cannot reach a state with current input
$k$.  Otherwise, the contradicting run emitted some output $x$ at $k$ whose
row the \commitment removed.  The \commitment leaves no row at $k$ that
\outrel $x$, whereas every row at $k$ after the emission does, so no row there
can be common to both runs.  \fullproof
\end{proof}

The proposition rules out reconciliation at the conflict input, not forever.  For example,
suppose $i\InputOrd j$, $\fiber{i}=\{o,x\}$ with $o$ and $x$ incomparable, and
$\fiber{j}=\{u\}$ with $u$ \outrelpart both.  The states that emit $o$ and $x$
cannot be reconciled at $i$, although accepting $j$ reconciles them.  As
explained after Definition~\ref{def:consistency}, consistency nevertheless
counts the conflict at $i$ as a failure.

The proposition says what a coordination mechanism must accomplish: prevent a
\commitment and a run that contradicts it from occurring independently.  This
requires distributed coordination when those runs can occur at different
participants.  The definitions diagnose the requirement, not its
implementation.  Locks, agreement protocols, ownership, and sequencers are
familiar implementations; using one need not apply a \commitment.

Consider a shared-nothing database in which each key is owned and locked by
one participant~\cite{stonebraker1986sharednothing}.  Ownership can keep
coordination local.  Ordering two transactions on one key requires
coordination because the owner may order them oppositely in another run, but
no other participant can produce that run.  A transaction spanning two keys
is different: the owners can order the same pair oppositely, so they must use
distributed coordination to agree on one order.

Two-phase commit adds no further coordination requirement.  Each owner's
prepare reports the \commitment it has already made; the coordinator's
monotone decision only detects whether those \commitments agree, aborting when
they do not.  A timeout abort before the last vote applies another \commitment,
like closing the polls.  It requires local coordination when a single
coordinator makes it.  It requires distributed coordination when the
coordinator is replicated for fault tolerance, as in Paxos
Commit~\cite{gray2006consensus}.  \subsection{The Complete CAP Theorem}
\label{sec:complete-cap}

Brewer conjectured that a distributed system cannot be consistent, available,
and partition tolerant at once~\cite{brewer2000towards}; Gilbert and Lynch
proved it for a linearizable register~\cite{gilbert2002cap}.  Their theorem
concerns one guarantee, not which guarantees can stay available and
consistent through a partition.  Complete CAP answers that broader question
for any declared guarantee.

To state the question, we must separate the outputs that participants must
keep available from the outputs that the specification permits.  A participant
acts on its local history during a partition, while consistency compares the
participants' outputs at a common \indesc of their inputs.  Their common
\outnoun need not be emitted by any participant.  It may exist only as the
client's combined observation.

Consider a register initialized to $v_0$.  During a partition, participant
$p$ receives a write $w(v)$ and emits its acknowledgment, while participant
$q$ receives a concurrent read $r$ and returns $v_0$.  A client that later
collects both observes the write acknowledgment together with the read answer
$v_0$.  The register specification permits this combined output at the input
containing both operations: a linearization may put the concurrent read before
the write.  The combined output \outrel both local outputs and provides a proof
of their consistency, although no participant need ever emit it.

We therefore encode a guarantee as an \spectype together with an
explicit relation $E\subseteq P$: the rows whose outputs must remain locally
available.  Outputs in $P\setminus E$ need only remain permitted as possible
common \outnouns.  Complete CAP says that such a guarantee can stay available
and consistent through partitions exactly when the specification satisfies
Lifted Monotonicity for $E$: outputs selected by $E$ on different sides of a
partition must have a common \outnoun in the specification at every common
\indesc of their inputs.  When $E=P$, this is Lifted Monotonicity itself.  A
partition does not change which \commitments a specification needs; it removes
the ability to make them jointly, since separated participants cannot jointly
apply a \commitment that another participant could contradict.

\begin{definition}[Partitioned history]
\label{def:partitioned-history}
A \emph{partitioned history} divides the participants into two or more
blocks and contains no receive whose matching send lies in another
block; participants still receive environment input and emit outputs.
\end{definition}

We fix a system, with a strategy for each participant of every finite
participant set and every choice of initial inputs, each depending only on
local history, and an adversarial environment that may supply
any legal input event to any participant.

\begin{definition}[Explicit relation, Availability]
\label{def:partition-notions}
A continuation of a partitioned history is \emph{local to} a participant $p$
when the system produces it, it adds only history events of $p$'s block and no
input events, and none of its receives takes a message sent from another
block.  An \emph{explicit relation} $E\subseteq P$ is a set of rows
$(i,o)$ whose output $o$ must stay available at input $i$; write $E(i)$ for
its fiber $\{o:(i,o)\in E\}$.  The system has \emph{Availability for $E$}
when, in every partitioned history it produces, each participant $p$, ending
in state $(i,R)$, can emit each output in $R(i)\cap E(i)$ in some continuation
local to $p$.  \emph{Availability} is Availability for $P$.
\end{definition}

\begin{definition}[Partition Consistency]
\label{def:partition-consistency}
The system has \emph{Partition Consistency} when the final cut of every
partitioned history it produces is confluent at every common \indesc of the
participants' current inputs.  The continuation that confluence asks for need
not stay partitioned and need not be produced by the system.
\end{definition}

\begin{definition}[Lifted Monotonicity for $E$]
\label{def:lm-for-A}
$\Spec$ satisfies \emph{Lifted Monotonicity for $E$} when every finite set of rows of $E$ with
a common \indesc $j$ has a common \outnoun in $\fiber{j}$.
\end{definition}

With $E=P$, Lifted Monotonicity for $E$ is the common-\indesc condition
(Definition~\ref{def:common-expansion-condition}).

\begin{theorem}[Complete CAP]
\label{thm:complete-cap}
For an \spectype $\Spec$ and an explicit relation
$E\subseteq P$, the following are equivalent:
\begin{enumerate}
  \item some system has Availability for $E$ and Partition Consistency;
  \item $\Spec$ satisfies Lifted Monotonicity for $E$.
\end{enumerate}
\end{theorem}

\begin{proof}[Proof sketch]
Consider the system whose strategies emit each output of
$E$ by the canonical output event of Section~\ref{sec:complete-calm}.  It has
Availability for $E$, and Lifted Monotonicity for $E$ gives the outputs it emits a
common \outnoun at every common \indesc of their inputs.  Conversely, given one
participant per row of $E$, each in its own block, Availability lets each
emit its row's output, and Partition Consistency then supplies the common
\outnoun.  \fullproof
\end{proof}

\begin{corollary}
\label{cor:complete-cap-S}
Some system has Availability and Partition Consistency if and only if $\Spec$
satisfies Lifted Monotonicity.  When it does, a \commitment-free system has
both: one whose strategies may emit any legal output with no
\commitment.
\end{corollary}
\proofinfull

The system in the proof needs no coordination.  Under Lifted Monotonicity for
$E$, the canonical event for an output selected by $E$ empties nothing and
excludes only rows outside $E$, which the system never emits, so it is
coordination-free (Definition~\ref{def:coordination-free}).  For $E=P$ it needs no
\commitment.  When Lifted Monotonicity for $E$ fails, some
output of $E$ waits for another block, sacrificing Availability, or is
emitted anyway, sacrificing Partition Consistency.

\paragraph{Gilbert and Lynch.}
Their theorem is a case of Complete CAP.  Consider a register specification
whose inputs are sets of operations with their real-time precedence, and whose
outputs answer every operation of the input as some linearization
does~\cite{herlihy1990linearizability}.  Let $p$'s input contain only a write
of $v\neq v_0$, and $q$'s only a read.  The fiber at $p$'s input holds only the
write's acknowledgment; the fiber at $q$'s input, which contains no write,
holds only the answer $v_0$ to the read.  At the common \indesc where the write
precedes the read, every output in the fiber answers that same read with $v$.  That contradicts
the answer $v_0$; it is not a later read observing a newer value.  So no output
there \outrel $v_0$, and $\Spec$ does not satisfy Lifted Monotonicity.  By Complete CAP, then, no system keeps both outputs available and
consistent across a partition, which is Gilbert and Lynch's theorem.
The full paper takes precedence from real time, extends
Complete CAP to systems that choose their responses when responses are not
forced, and derives
the theorem in Gilbert and Lynch's asynchronous model~\cite{hellerstein2026completecalmfull}.

% ======================================================================
\section{Composition and Proper Coordination}
\label{sec:composition}

Coordination poses a puzzle.  A non-monotone program needs coordination, yet
coordination is itself non-monotone.  A key-value store with mutable keys is
non-monotone~\cite{alvaro2011bloom}.  State-machine
replication~\cite{schneider1990smr} wraps it in a log.  Clients send requests,
the log puts the requested commands in one order, and the key-value store
executes the log and answers each request.  The log's specification is also not
monotone: when two commands arrive concurrently, either may be appended next,
and the two alternative orders have no common \outnoun.  Yet replicated stores are
consistent, coordinating only in the log.  Multi-Paxos agrees on each
log entry~\cite{lamport2001paxos}, and the replicas then execute the log
with no further agreement.

That puzzle is explained by our framework.  Specifications are
tables, so they compose by equijoin, and composition preserves Lifted
Monotonicity.  Because \commitments are explicit, we
can also ask which component makes a composition's \commitments.  When the
upstream component makes every \commitment the composition needs, and the
downstream needs none on the inputs the upstream offers, we call the
coordination \emph{proper}
(\S\ref{sec:proper-coordination}).
\label{sec:sequential-composition}

Intuitively, sequential composition feeds the output of one specification to
the input of the next.  The composite is again a table: it pairs an input of
the first specification with every output that the second permits for some
intermediate value that the first permits.  Formally, let
$\Spec_1=(I,\InputOrd,M,\preceq_M,P_1)$ and
$\Spec_2=(M,\preceq_M,O,\OutputOrd,P_2)$ be \spectype[s]
sharing the ordered domain $(M,\preceq_M)$.  Their \emph{sequential
composition} has table
\begin{equation}
  P_2\circ P_1
  =\pi_{I,O}(P_1\Join_M P_2)
  =\{(i,o):\exists m.\ (i,m)\in P_1\land(m,o)\in P_2\},
  \label{eq:spec-composition}
\end{equation}
where $\pi$ and $\Join$ are relational projection and equijoin.  Since $P_1$
is total on $I$ and $P_2$ is total on $M$, every $i\in I$ has a nonempty fiber
in $P_2\circ P_1$.  When the two stages
run as a pipeline, $\Spec_2$ takes the outputs of $\Spec_1$ as its inputs, and
we say that $\Spec_1$ \emph{guards} $\Spec_2$.

\begin{lemma}[Composition]
\label{lem:composition}
If $\Spec_1$ and $\Spec_2$ satisfy Lifted Monotonicity, then so does the
\spectype with table $P_2\circ P_1$.
\end{lemma}

\begin{proof}[Proof sketch]
Take rows of $P_2\circ P_1$ whose inputs have a common \indesc $j$.  The
common-\indesc condition for $\Spec_1$ gives an intermediate value at $j$
that \outrel each row's, and that value is a common \indesc
at which the condition for $\Spec_2$ supplies the common \outnoun.  \fullproof
\end{proof}

By associativity the lemma extends to finite pipelines.

\subsection{Composite CALM}
\label{sec:composite-calm}

By Complete CALM, $P_2\circ P_1$ is consistent exactly when it
satisfies Lifted Monotonicity, and then no output needs a \commitment.
That test treats the composite as a single specification.  Suppose instead that
each stage runs on its own, with its runs taking as input the outputs of the
stage before.  A stage may then need a \commitment even when the composite
needs none.  A composite can be jointly monotone, with no stage needing a
\commitment, in two ways.  First, if every stage is monotone,
as when a monotone client query observes CRDT state
(\S\ref{sec:crdt-instance}), then Lemma~\ref{lem:composition} applies.
Second, the upstream can filter out the inputs on which the downstream is not
monotone: a counter capped at $9$ never produces the counts at which a test
of whether the count is below $10$ answers no.  These cases do not cover every monotone composite
(see \extended).  When none applies, some stage may have to issue a \commitment even though the
composite is monotone.

\subsection{Proper Coordination}
\label{sec:proper-coordination}

The replicated store of the opening extends the filtering case.  The
key-value store is not monotone
on its own, but it is monotone on what the log offers: a log fixes the order
of commands, and a longer log only extends the responses.  Unlike the
filtering case, the upstream must issue a \commitment to offer only such inputs.

Let $\Spec_1$, which guards $\Spec_2$, be run by a system with any basis, and
let $\Spec_2$ run under a system of its own.  Write $\Spec_2|_{\Spec_1}$ for $\Spec_2$ restricted to the outputs of
$\Spec_1$.

\begin{definition}[Proper Coordination]
\label{def:proper-coordination}
A pipeline in which $\Spec_1$ guards $\Spec_2$ is \emph{proper} when
$\Spec_2|_{\Spec_1}$ needs no \commitment: every output at every state it
reaches without \commitment is legal.
\end{definition}

\begin{theorem}[Proper Coordination]
\label{thm:proper-coordination}
Fix any system for $\Spec_1$.
\begin{enumerate}
  \item The pipeline is proper exactly when $\Spec_2|_{\Spec_1}$ satisfies
        Lifted Monotonicity.
  \item If it is proper, $\Spec_2|_{\Spec_1}$ under the empty basis follows
        every run of $\Spec_1$.  The result is a run of $P_2\circ P_1$, and
        whenever runs of $\Spec_1$ can be brought to a common state, the
        composite runs are confluent.
  \item If it is proper and $P_2\circ P_1$ violates Lifted Monotonicity,
        $\Spec_1$ cannot emit some output without a \commitment.
\end{enumerate}
\end{theorem}

\begin{proof}[Proof sketch]
(1) is Complete CALM and the output-legality boundary
(Theorem~\ref{thm:complete-calm}, Proposition~\ref{prop:exposure-boundary})
applied to $\Spec_2|_{\Spec_1}$.  (2) Each output event of $\Spec_1$ keeps
only rows above its output, so each input $\Spec_2$ receives lies above the
previous one, and by (1) every answer is legal.  (3) The output-legality boundary
of the composite fails at a canonical state, which is the image of a canonical
state of $\Spec_1$.  \fullproof
\end{proof}

The theorem holds whatever basis $\Spec_1$ uses, and a proper pipeline is
itself a guard for a further stage.  Locks in
the key-value store would be pure overhead: by (1), the store restricted to
logs needs no \commitment.  The full paper extends these results to more complex
compositions than sequential stages, including
stages with inputs of their own and arbitrary acyclic graphs of stages, checked node
by node much as types are~\cite{hellerstein2026completecalmfull}.

% ======================================================================
\section{Related Work and Mappings}
\label{sec:instances}
\label{sec:related}

Complete CALM provides a unified lens for many results across literatures.
For the key related ideas, we provide mappings here to our specification model.
Each setting declares an input, an output, and their orders before
applying either CALM theorem.  The full paper~\cite{hellerstein2026completecalmfull}
gives each mapping in detail, with deployments, response
obligations, and \commitment bases for the transaction protocols 2PL and MVTO,
and applications to distributed-systems protocols such as registers, $k$-set agreement, renaming, and snapshots.

\subsection{Query monotonicity}
\label{sec:query-instance}

As previously discussed (\S\ref{sec:functional-calm}),
the CALM Theorem of Ameloot et al.\ proves an operational result for
relational languages on transducer networks: coordination-free evaluation is possible
iff a program is monotone,
under its data-placement and execution assumptions~\cite{ameloot2013relational,ameloot2015weaker}.
Functional CALM does not replace the details of that theorem.  It abstracts
the result to a property that can be compared with specifications outside
Datalog. Subsequent extensions to the original CALM Theorem become changes to the input domain
and its order.
When transducers know
how data is placed, more queries become coordination-free, among them the
non-monotone win-move query~\cite{zinn2012winmove}, and Ameloot et al.\
characterized the new classes by domain-distinct and domain-disjoint
monotonicity~\cite{ameloot2015weaker,zinn2012weak}.  Each is Functional CALM
under an input order in which $I\cup J$ \inrel $I$ only for the $J$ that the
placement policy leaves a node unable to rule out.  Functional CALM takes the
order as given.

Baccaert and Ketsman's \emph{behaviors} map located input facts to sets of
permitted located outputs~\cite{baccaert2023dcbq,baccaert2026spectrum}; under
inclusion, a behavior is an \spectype.  Their characterization,
\emph{node monotonicity}, is the common-\indesc condition for two rows at
disjoint networks, and Complete CAP extends it to overlapping networks
(see \extended).

\subsection{CRDTs: monotone queries over convergent state}
\label{sec:crdt-instance}

A CRDT guarantees that replicas receiving the same updates converge to
equivalent states, but the Strong Eventual Consistency property it offers
provides no guarantee on what the environment may observe.  Laddad et al.\
proposed monotonicity as the right test for safe observation from CRDT
state~\cite{laddad2022keep}.  We make the proposed test precise with
Functional CALM.  Let $(D,\sqsubseteq,\sqcup)$ be the state join-semilattice.
The input is the finite set $U$ of immutable payloads received so far, ordered
by inclusion, and the replica's state is $s(U)=s_0\sqcup\bigsqcup U$.  A
client observation $q:D\to(Y,\preceq_Y)$ gives the client-visible
specification $q\circ s$.  By Functional CALM, $q\circ s$ is consistent
exactly when $U\subseteq V$ implies $q(s(U))\preceq_Y q(s(V))$.  Monotonicity
of $q$ on all of $D$ suffices, but $q$ need only be monotone between reachable
states.  The query $\mathit{Add}\setminus\mathit{Del}$ over a two-phase set
fails the test: a deleted item grows $U$ and shrinks the answer.  CRDT states
form a join-semilattice, but $q$ needs less: Complete CALM requires
only \emph{some} common \outnoun of query outputs, not a least one.

\subsection{Highly available transactions}
\label{sec:hat-instance}

We take the assumptions of Bailis et al.\ as given~\cite{bailis2013hat}:
items are partially replicated; a client sends each operation to some replica
of its item; a transaction's writes carry a unique client timestamp, last
writer wins; the system never aborts on its own; and a replica answers a read
with the highest-timestamped committed write of the item it has received.

\begin{definition}[HAT specification]
\label{def:hat-spec}
An input is a history ordered by happens-before: reads tagged with their
replica, commits with their timestamp and write set, and deliveries of
committed writes to replicas.  Inputs are ordered by
downward-closed extension, and outputs are sets of read answers, ordered by
inclusion.  For an isolation level $L$, defined by the phenomena it
proscribes~\cite{adya1999weakconsistency}, let $\Spec_L$ have table $P_L$,
which permits an output when it
extends to a history that answers every read, orders versions by timestamp,
and exhibits none of them.  The explicit relation $E$ holds, at each
input, the answers the read rule gives.
\end{definition}

\begin{theorem}[HAT levels]
\label{thm:hat-levels}
Suppose each transaction's timestamp exceeds that of every version it reads.
For read uncommitted and read committed, $E\subseteq P_L$ and $\Spec_L$
satisfies Lifted Monotonicity for $E$, so by Complete CAP these levels stay
available and consistent through partitions.  For any level that proscribes
Lost Update, including snapshot isolation, repeatable read, and
serializability, no explicit relation that answers every read satisfies
Lifted Monotonicity for it.
\end{theorem}
\proofinfull

The condition is needed: with timestamps fixed at transaction
begin, as in the pseudocode of~\cite{bailis2013hatextended}, a transaction
that reads a higher-timestamped version and overwrites it forms a G1c
cycle.

Mahajan et al.\ showed that no consistency semantics stronger than real-time
causal consistency can be enforced in an always-available, one-way convergent
system~\cite{mahajan2011consistency}.  With real-time precedence recorded in
the input, both halves of their result follow from Complete CAP (\extended).

\subsection{Invariant confluence}
\label{sec:iconfluence-instance}

In the model of Bailis et al.~\cite{bailis2014coordination}, a database state
is a set of versions, replicas merge states by union, and a set $T$ of
transactions is \emph{invariant confluent} for an invariant $\mathcal I$ when
the merge of any two reachable states satisfies $\mathcal I$.  We construct a
specification $\Spec_{\mathcal I}$ whose inputs are locally valid
\emph{transaction histories}, partial orders of transactions in which one
precedes another when the second's replica saw the first, ordered
by extension.  Outputs are database states, ordered by inclusion;
$P_{\mathcal I}$ permits at $h$ the state of any downward-closed
sub-history of $h$ whose merges all satisfy $\mathcal I$.

\begin{proposition}[Invariant confluence and Lifted Monotonicity]
\label{prop:iconfluence-lm}
$T$ is invariant confluent if and only if $\Spec_{\mathcal I}$ satisfies Lifted
Monotonicity.
\end{proposition}
\proofinfull

In Bailis et al.'s payroll example, two
replicas each assign ID~7 to a new employee; the history in which neither
creation saw the other \inrel both one-creation histories, but no state it
permits \outrel both.

\subsection{Operational models and scope}
\label{sec:scope}

Automata are a traditional way to model operational behavior
abstractly~\cite{lynch1987io,lynch1996distributed}.  Our model keeps only the
inputs and outputs of that behavior and drops the operational details.  The
mapping is straightforward: for each input, record the outputs of the runs that
consume it.  Two standard approaches then apply, both used widely and with
success.  Model checkers explore runs on bounded
inputs~\cite{flanagan2005dpor,newcombe2015aws};
their input/output pairs become a finite specification in our model that can be tested for
violations of Lifted Monotonicity.  The only additional requirement is on the
checker's bounds: any input it reports on must have had all its executions
explored. The more comprehensive approach is to use
formal proofs to declare a specification and show, by the standard obligations
of refinement~\cite{abadi1991refinement,lynch1995forward}, that it represents
the machine: every execution's output is dominated by a declared one, and every
declared row is produced by some execution.  Such proofs are routinely
mechanized~\cite{hawblitzel2015ironfleet}, and Lifted Monotonicity is then
checked on the declared specification; \extended{} gives the details.  Work that
models systems as state machines therefore maps naturally into our model, where
Complete CALM assesses its consistency directly.

  Our conditions are semantic.  Confluence guarantees that
a common \outnoun exists, not that a deployed system reaches it; protocols,
liveness, and cost belong to the operational model in which a specification
is realized.  Free termination likewise corresponds to a determined canonical
state in our model, but Power et al.\ additionally study fairness and eventual
response, liveness concerns whose semantic generalization remains future
work~\cite{power2025freetermination}.

% ======================================================================
\section{Discussion}
\label{sec:discussion}

The computing literature is rich with coordination mechanisms for managing
concurrency within and across machines. The CALM Conjecture asked a question
still latent at the time: when are these mechanisms required, and why? It drew
on declarative networking~\cite{loo2009declarative}, which, in Codd's tradition
of separating specification from implementation~\cite{codd1970relational}, wrote
the desired effects of protocols as queries to be analyzed and optimized into a
variety of mechanisms. Could a signature property of declarative
languages---monotonicity---be the criterion?

The intuition was sharpened by a related friction in the literature on parallel
query processing. The Google MapReduce
paper~\cite{dean2004mapreduce} departed from prior work on parallel query
processing~\cite{wilschut1991dataflow,graefe1990volcano}, conflating data
redistribution (shuffle) with barrier coordination required for a non-monotone
reduce operator. This coupling of unrelated concerns required monotone operators
like join to block in the MapReduce model unnecessarily.  Cost models that count
rounds, from BSP supersteps to MapReduce rounds, carry the same
conflation~\cite{valiant1990bsp,karloff2010mrc}.

The CALM Conjecture hypothesized that the database community's more judicious
use of barriers could be generalized: coordination---of any kind---is required
if and only if the program is non-monotone. Ameloot et al.\
proved it, but only for relational languages, which to date have
still not become mainstream beyond query processing (more on this below).

This paper moves the declarations one level higher, from queries to views: all
that matters about a specification is the single view mapping a program's inputs
to its output. Eliding the query logic frees the analysis from any particular
language and broadens what can be specified. Prior work focused on implementing
complex programs
declaratively~\cite{loo2005overlays,loo2009declarative,chu2007declarative,alvaro2011bloom}.
This was interesting ``declarative engineering'' with concrete benefits in
analysis and optimization. But the higher-level model here also captures abstract
families of guarantees, like CAP and isolation levels, which have no application or
protocol in mind.

Lifting the abstraction improves the power of analysis as well. Prior program
analyses became enmeshed in the details of implementations, however
declarative. Bloom~\cite{alvaro2011bloom} flagged syntactic non-monotonicity in
programs and conservatively injected coordination code in front of it.  Yet
coordination code is itself non-monotone, so adding it to a Bloom program only
adds more non-monotone syntax.  Subsequent work moved the question from syntax to declared contracts and
invariants~\cite{sivaramakrishnan2015quelea,kaki2018alone}.  Their tools pick,
for each operation or transaction, the weakest consistency or isolation level
that the store offers and that still satisfies the declared properties.  We
ask what the specification itself requires, independent of any menu of
levels.  By Canonical Run Universality (Lemma~\ref{lem:canonical-equivalent}), the answer holds for every mechanism.  Our
composition results need no program logic, and the criteria are exact on
specifications, not conservative.

A limitation of prior CALM work was that it was constructive on only one side of
its boundary.  Monotone programs need no coordination; for a non-monotone
program, CALM said only that coordination was required somewhere, not what it
had to accomplish or where it belonged.  Complete CALM says more.  A
specification that is not monotone needs specific \commitments: the \commitments
its outputs require. Every coordination mechanism, from locks to
consensus, is a way of applying these \commitments.  The same analysis says which outputs
need \commitments, and whether another participant could contradict them, so
that agreement is required and not just a local decision.  Proper coordination
then places the \commitments in a composition: it identifies the stage that
must make them, as the log does in state-machine replication, and certifies
that the stages behind it need none of their own.

Underlying these issues is a reassessment of
tradition: what is a declarative specification, and how little does it require
to be useful? The question is timely. Coding agents increasingly write
imperative code, often rife with bugs, which has renewed interest in formal
specification languages like Lean~\cite{demoura2021lean4}. A higher-level
language keeps agents out of the incidental complexity of implementations and
pairs them with proofs of correctness: neurosymbolic
computing~\cite{an2025neurosymbolic}. This is a renewed opportunity for logic,
dataflow, and functional languages. Today these languages typically specify
protocols and programs in detail, as in declarative networking; this paper
points to specification at a higher level still, opening analyses like ours and
potential program synthesis and optimization.

% ======================================================================
\section{Conclusion}
\label{sec:conclusion}

A specification abstracts execution details away, and the results here reduce
to a single question: can a participant act on what it knows without being
contradicted by what it learns later?  Mutexes, locks, atomic instructions, and
agreement protocols prevent such contradictions by issuing \commitments.
Functional CALM answers the question for functions: a functional specification
stays confluent as input arrives exactly when it is monotone, and then no
output needs a \commitment.  This generalizes the classic result from logical fact
sets to arbitrary input and output orders.

Complete CALM extends the answer to specifications that offer several outputs
for one input.  The test is Lifted Monotonicity, which combines Relational
Monotonicity with Directed Outputs.  When it fails, some reachable output
requires a \commitment.  At each input it changes, a \commitment either empties
the fiber or removes particular output rows.  Because \commitments are explicit, the model
can also check where a composition coordinates
(\S\ref{sec:composition}).  A totally ordered log and a
key-value store are both non-monotone at their ordinary specifications, yet
when the log guards the store, the log's \commitments are all the composition
needs and the store makes none.  Multi-Paxos is one way to distribute those
\commitments.

Complete CAP is the partition form of the same equivalence.  A system can keep
the outputs selected by a guarantee's explicit relation available, and
consistent, through partitions exactly when they have a common \outnoun at
every common \indesc of their inputs; when the relation selects every output,
this is Lifted Monotonicity.

The generality of our specification model opens a range of new questions
regarding the nondeterminism of specifications and their \commitments. One direction
extends the boolean criteria of this paper to metrics: a Determination
Complexity theory that assesses \emph{how many} \commitments are intrinsically
needed to determine a specification. Another is Determination Provenance: a
given output---or its absence---is determined not only by inputs, but also by
the \commitments made in a run, which brings formal models of provenance to
concurrent systems. A third direction is to provide general-purpose static
analysis of CALM and proper coordination in a real programming framework.

\begin{acks}
Thanks to Conor Power, whose persistent questioning of the CALM canon prompted this work.
Much gratitude to Peter Alvaro, Natacha Crooks, Chris Douglas, Tyler Hou, Martin
Kleppmann, Paris Koutris, Edward Lee, Shulu Li, Vincent Liu, David Chu McElroy,
Sarah Morin, and Dan Suciu for helpful feedback. Thanks also to
anonymous reviewers of this and previous submissions for frank feedback that pushed
this work deeper.

Generative AI tools (Infinity, Kiro, Claude and ChatGPT) were used as interactive writing and
reviewing assistants during the preparation of this work, including drafting
and revising text, and checking mathematical arguments with Lean proofs.  All
content was reviewed, validated, and approved by the author.
\end{acks}

\bibliographystyle{ACM-Reference-Format}
\bibliography{references,cidr-extra}

\end{document}
